\documentclass[11pt]{article}
\usepackage[margin=0.85in]{geometry}
\usepackage{amsmath,hyperref}
\title{Comparing fixed-OH and joint methane--OH inversions}
\author{Reading note on Feng et al. (2023)}\date{}
\begin{document}\maketitle
\noindent\textit{Methane emissions are predominantly responsible for
record-breaking atmospheric methane growth rates in 2020 and 2021.}
Atmospheric Chemistry and Physics 23, 4863--4880.
\href{https://doi.org/10.5194/acp-23-4863-2023}{doi:10.5194/acp-23-4863-2023}.
Appendices A--C are included in the published article; no separate
supplement is listed.

\paragraph{Method.}
The ensemble Kalman filter jointly assimilates Leicester GOSAT v9
CH$_4$:CO$_2$ column ratios and surface measurements of both gases.
It therefore avoids assigning the proxy ratio to methane using prescribed
CO$_2$ alone. GEOS-Chem uses MERRA2 meteorology at
$2^\circ\times2.5^\circ$. Monthly emission factors cover 476 land
and 11 ocean regions, with a four-month assimilation window.
Prior methane flux uncertainty is 60\%, with 300-km spatial and one-month
temporal error correlation. Alongside a fixed-OH control, a joint inversion
for 2018--2021 optimizes six monthly zonal OH factors with 3\% prior
uncertainty. Polar OH is not adjusted. Wetlands and rice follow WetCHARTs;
other anthropogenic priors use EDGAR v4.3.2 for 2012.

\paragraph{Findings.}
Table~1 gives fixed-OH emissions of 583.7, 610.7, and 604.5 Tg yr$^{-1}$
for 2019, 2020, and 2021, respectively. The joint inversion gives
585.3, 603.2, and 603.7 Tg yr$^{-1}$, so its 2019-relative anomalies
are 17.9 and 18.4 Tg yr$^{-1}$. Published OH changes relative to the
prior are $+0.91$, $-0.52$, and $+0.62$\%; the text reports changes
of $-1.43$ and $-0.29$ percentage points relative to 2019. The
34\% attribution in 2020 is the reduction of the inferred emission
anomaly from 27.0 to 17.9 Tg yr$^{-1}$ when OH is optimized. This
is a comparison of two inversions, not a direct evaluation of
$-m_0\Delta k_{\mathrm{OH}}$. Their posterior regional flux--OH
error correlations are generally below 0.05 in magnitude.

Eastern Africa contributes about 14 Tg yr$^{-1}$ in the fixed-OH
2020 comparison. China and India decrease by about 6 and 2 Tg yr$^{-1}$.
Tropical South America and temperate North America matter more in
2021. Regional water-storage and emission anomalies correlate, but
the association alone does not distinguish microbial sectors.

\paragraph{Comparison decisions.}
Appendix~B explicitly distinguishes growth calculated between calendar-year
endpoints from differences between annual means: the latter shift the
effective timing by roughly half a year. We retain this distinction and
use each inversion's own 2019 emissions. Annual posterior errors of
roughly 11--13 Tg yr$^{-1}$ must not be treated as uncertainty of
year-to-year differences without the cross-year covariance. Table~1
can be transcribed exactly; no gridded posterior download is linked.
The 9.1 Tg yr$^{-1}$ fixed-versus-joint difference in 2020 is retained
as a sensitivity diagnostic and is not labeled a measured OH forcing.
TCCON validation and the sensitivity experiments support the inferred
large-scale flux patterns, conditional on the assumed OH parameterization.
\end{document}
